\part{The Firm: Production and Cost}

\chapter{Production and Technology}
\label{chap:production}

\section{Learning objectives}

After completing this chapter, you should be able to:
\begin{itemize}
  \item define production functions, inputs, and technology;
  \item calculate marginal and average products;
  \item interpret isoquants and the marginal rate of technical substitution;
  \item distinguish the short run from the long run; and
  \item explain returns to scale.
\end{itemize}

\section{Firms and production}

A \keyterm{firm} transforms inputs into outputs. A baker combines labour, ovens, ingredients, and knowledge to produce bread; a software firm combines programmers, computers, data, and organisational knowledge to produce digital services. \keyterm{Technology} describes the set of methods available for transforming inputs into outputs.

A \keyterm{production function} gives the maximum output that can be produced from a set of inputs using a given technology. With labour $L$ and capital $K$,
\[
q=f(L,K).
\]
The production function is a technological relationship, not a statement about demand or prices. It identifies what is feasible; firms later choose among feasible methods by considering costs and revenue.

\section{Marginal and average product}

In the short run, at least one input is fixed. Suppose capital is fixed and labour varies. The \keyterm{marginal product of labour} is the additional output from one more unit of labour:
\[
MP_L=\frac{\partial f(L,K)}{\partial L}.
\]
The \keyterm{average product of labour} is output per unit of labour:
\[
AP_L=\frac{q}{L}.
\]

Initially, adding workers may raise marginal product through specialisation. Eventually, if capital remains fixed, workers may crowd one another and marginal product falls. This is the \keyterm{law of diminishing marginal product}: holding other inputs fixed, the marginal product of a variable input eventually decreases as more of it is used. It is not a claim that total output must fall; total output can continue to rise while doing so more slowly.

\section{Isoquants}

An \keyterm{isoquant} shows all combinations of inputs that produce the same output. It is the production-side analogue of an indifference curve. A firm on an isoquant can substitute labour for capital while maintaining output.

With monotonic technology, isoquants slope downward: if the firm uses less capital, it needs more labour to produce the same output. Higher isoquants represent higher output. Isoquants do not cross, because the same input combination cannot produce two different maximum output levels with a fixed technology.

\section{The marginal rate of technical substitution}

The \keyterm{marginal rate of technical substitution} of labour for capital is the amount of capital that can be reduced when one more unit of labour is employed without changing output. Its absolute value is the slope of an isoquant:
\[
MRTS_{L,K}=\frac{MP_L}{MP_K}.
\]

For many technologies, the MRTS diminishes as labour replaces capital. When a production process already uses much labour and little capital, an extra worker cannot easily replace a machine. This produces convex isoquants.

\section{Short run and long run}

The \keyterm{short run} is a period in which at least one input is fixed. The \keyterm{long run} is a period long enough for all inputs to be varied. These are economic concepts, not fixed calendar periods. A restaurant may adjust staff within days but take years to change location; a digital firm may scale server capacity much more quickly.

The distinction matters because firms have more flexibility in the long run. It also explains why a decision that appears inefficient in the short run may be sensible when some costly commitments cannot yet be changed.

\section{Returns to scale}

Returns to scale describe what happens when all inputs change proportionally. A production function has \keyterm{constant returns to scale} if doubling every input doubles output, \keyterm{increasing returns to scale} if output more than doubles, and \keyterm{decreasing returns to scale} if output less than doubles.

For the Cobb--Douglas function $q=AL^\alpha K^\beta$, returns to scale depend on $\alpha+\beta$. They are increasing if this sum is greater than one, constant if it equals one, and decreasing if it is less than one.

Increasing returns may arise from specialisation, indivisible fixed costs, or learning. They are important for understanding large firms and concentrated industries. They do not contradict diminishing marginal product: the first compares proportional changes in all inputs, while the second holds at least one input fixed.

\section{Common technologies}

The Cobb--Douglas production function $q=AL^\alpha K^\beta$ allows smooth substitution between inputs. The Leontief production function $q=\min\{aL,bK\}$ describes fixed-proportion production: additional workers without the necessary machines do not increase output. These forms are simplifications, but they make the relevant technological trade-offs clear.

\section*{Chapter Review}

\subsection*{Key concepts}
\begin{itemize}
  \item \textbf{Production function}: the maximum feasible output from inputs.
  \item \textbf{Marginal product}: additional output from an additional input.
  \item \textbf{Average product}: output per unit of input.
  \item \textbf{Isoquant}: input combinations producing the same output.
  \item \textbf{MRTS}: technological trade-off between inputs along an isoquant.
  \item \textbf{Returns to scale}: output response to proportional changes in all inputs.
\end{itemize}

\section*{Exercises}

\subsection*{Conceptual questions}
\begin{enumerate}[label=\arabic*.]
  \item Define a production function. Why is it not a demand curve?
  \item Distinguish marginal product from average product.
  \item Explain the law of diminishing marginal product.
  \item Distinguish diminishing marginal product from decreasing returns to scale.
  \item Explain why isoquants normally do not cross.
\end{enumerate}

\subsection*{Graphical exercises}
\begin{enumerate}[resume,label=\arabic*.]
  \item Draw three convex isoquants for increasing output and label labour and capital axes.
  \item Draw L-shaped isoquants for a Leontief technology. Explain why substitution is impossible away from the kink.
  \item Sketch total product, marginal product, and average product curves consistent with diminishing marginal product.
\end{enumerate}

\subsection*{Analytical exercises}
\begin{enumerate}[resume,label=\arabic*.]
  \item Let $q=2L^{1/2}K^{1/2}$. Find $MP_L$, $MP_K$, and $MRTS_{L,K}$.
  \item For $q=3L^{0.6}K^{0.7}$, determine returns to scale.
  \item A firm has $q=\min\{2L,K\}$. What combinations of $L$ and $K$ produce $q=20$ efficiently?
\end{enumerate}

\subsection*{Exam-style multiple-choice questions}
\begin{enumerate}[resume,label=\arabic*.]
  \item Diminishing marginal product occurs when: (a) all inputs rise proportionally; (b) the marginal product of a variable input falls while other inputs are fixed; (c) total output immediately falls; or (d) prices increase.
  \item An isoquant shows: (a) equally costly input combinations; (b) equally preferred consumption bundles; (c) input combinations producing the same output; or (d) market-equilibrium prices.
  \item If doubling all inputs more than doubles output, the technology has: (a) decreasing returns to scale; (b) constant returns to scale; (c) increasing returns to scale; or (d) diminishing marginal utility.
\end{enumerate}

\subsection*{PPE discussion question}

\noindent Large-scale production can lower costs but may also concentrate economic power. Discuss whether increasing returns to scale provide an argument for regulation, public ownership, or neither.
